\section{Capacity of deterministic shared decoders}\label{sec:capacity}

Throughout this section $F$ is a field, $d,b\ge1$, and
$V=(F^b)^{(\Z/d)\times(\Z/d)}$ has dimension $N=bd^2$.
A partition $p$ of the receiver indices has $M$ nonempty classes.
The source $C\le V$ is linear, of dimension $D$; the shared decoders
$f_J:V\to(F^b)^d$ are arbitrary total functions satisfying
$f_{p(c)}(T^cX)=\operatorname{row}_cX$ for $X\in C$.
All lanes may be coupled. This distinction between a linear source and a
linear decoder will remain explicit.

\subsection{Intersect inside the covariance space}

For a nonempty set $E\subseteq\Z/d$, put
\[
 C_E=\bigcap_{e\in E}T^{-e}C.
\]
Join distinct indices $c,c'$ when $p(c)=p(c')$ and $c'-c\in E-E$;
write $\chi(E)$ for the number of components. Let $H_E$ be the subspace
defined by
\[
 \operatorname{row}_{c'}X=\sigma^{c'-c}\operatorname{row}_cX
 \quad\text{on every edge}.
\]

\begin{lemma}[Deterministic translate pigeonhole]\label{lem:deterministic-tp}
\formal{CapacityAmbient.lean}{IsSharedCode.component_finrank_le, finrank_covarianceSpace, covarianceSpace_rot}
For every nonempty $E$, $C_E\subseteq H_E$, and
$\dim H_E=bd\chi(E)$. Consequently
\[
 |E|D\le (|E|-1)N+bd\chi(E).
\]
Every $H_E$ is invariant under every column rotation $T^t$.
\end{lemma}
\begin{outline}
If $c+e=c'+e'$ with $e,e'\in E$ and $p(c)=p(c')$, apply the two
correct decoder equations to $T^eX,T^{e'}X\in C$. Their arguments are
the same full matrix, so
$\sigma^e\operatorname{row}_cX=\sigma^{e'}\operatorname{row}_{c'}X$.
This gives the edge equation without using linearity of the decoder.
After replacing row $c$ by $\sigma^{-c}\operatorname{row}_cX$, the
equations say simply that normalized rows are constant on each component.
One arbitrary vector of dimension $bd$ per component gives all solutions;
there is no further condition around a cycle. This proves the dimension.
Codimension is subadditive under intersection, and each $T^e$ is invertible,
so $\dim C_E\ge |E|D-(|E|-1)N$. Finally, applying the same rotation
$\sigma^t$ to all rows preserves every edge equation because the rotations
commute. This proves invariance.
\end{outline}

The covariance argument is that of the translate-pigeonhole theorem in
\cite{SeventeenthPaper2026}. Its use of decoder values does not require those
decoders to be linear. The smaller invariant ambient in the next theorem
is the additional resource.

\begin{theorem}[Iterated translate bound]\label{thm:iterated-tp}
\formal{CapacityIteration.lean}{IsSharedCode.iterated_finrank_le, componentCount_zero}
Choose shifts $t_1,\ldots,t_k$, allowing repetitions, and set
\[
 E_0=\{0\},\qquad E_j=E_{j-1}\cup(E_{j-1}+t_j),\qquad
 n_j=bd\chi(E_j).
\]
Then $n_0=N$ and
\begin{equation}\label{eq:iterated-tp}
 D\le\sum_{j=0}^{k-1}\frac{n_j}{2^{j+1}}+\frac{n_k}{2^k}.
\end{equation}
\end{theorem}
\begin{outline}
Write $C_j=C_{E_j}$. The set identity gives
$C_j=C_{j-1}\cap T^{-t_j}C_{j-1}$. Both spaces on the right lie in
$H_{E_{j-1}}$, by the preceding lemma and rotation invariance. They have
the same dimension. Grassmann's formula inside this ambient gives
$2\dim C_{j-1}\le n_{j-1}+\dim C_j$. Iterate these inequalities and
use $\dim C_k\le n_k$. Repeated subset sums are counted once in $E_j$;
their collisions do not affect either the set identity or the argument.
\end{outline}

The theorem bounds actual simultaneous intersections. It assumes neither
independence of the constraints nor correctness of $H_E$ itself. It can
be combined with every applicable direct translate bound.

\subsection{A uniform bound for two classes}

Suppose now that $p$ has at most two labels. For $s\in\Z/d$, let
$g_s=\gcd(d,s)$, where $g_0=d$, and
$Q_s=|\{c:p(c)\ne p(c+s)\}|$.
On the cycles of addition by $s$, let $R_j(s)$ count proper monochromatic
runs of length at least $j$. A proper run is preceded by a color change;
a monochromatic cycle contributes no runs. Thus $R_1(s)=Q_s$.

\begin{lemma}[Run components and doubled steps]\label{lem:binary-runs}
\formal{CapacityRunsComponents.lean}{componentCount_le_runs_add_two_gcd, componentCount_le_changes_add_gcd}
\formal{CapacityRuns.lean}{CapacityRuns.changeCount_double, CapacityRuns.weighted_run_budget}
For every integer $r\ge1$, the same-color graph with edges at $s$-cycle distance at most $r$ has
at most $R_r(s)+2g_s$ components. At radius one it has at most $Q_s+g_s$.
For every $s$, including zero and steps with short cycles,
\begin{align}
 Q_{2s}&=2R_2(s),\label{eq:doubled-runs}\\
 Q_s+R_2(s)+2R_4(s)+2R_{13}(s)&\le d.\label{eq:run-budget}
\end{align}
\end{lemma}
\begin{outline}
Work on one $s$-cycle. If its length is at most $2r$, the radius-$r$
graph is complete within each color and has at most two components.
Otherwise list the positions of one occurring color cyclically. Consecutive
positions are connected when their intervening gap is at most $r$.
Every other short edge is a path of these consecutive short edges, so it
adds no component merger. If $B>0$ gaps are longer than $r$, there are
$B$ components; if no gap is long, there is one. A singleton has its
one long cyclic gap and one component. For a nonconstant binary cycle,
these long gaps are exactly the opposite-color runs of length at least
$r$. Summing the two colors proves the first estimate. At radius one a
nonconstant cycle has exactly one component per run; a monochromatic
cycle adds one component and no change. Summing cycles proves both claims.

For the identity, write $z_c=1_{p(c)\ne p(c+s)}$.
Then $1_{p(c)\ne p(c+2s)}=z_c\mathbin{\mathrm{xor}}z_{c+s}$.
In a nonconstant, nonalternating cycle, the changes of $z$ number twice
its nonempty runs of ones, each ending at the start of a $p$-run of length
at least two. Constant and alternating cycles give zero on both sides.
This also handles cycles of lengths one and two. Finally a proper run
of length $L$ contributes
$1+1_{L\ge2}+2\,1_{L\ge4}+2\,1_{L\ge13}$ to the left of
\eqref{eq:run-budget}. Its values are $1,2,4,6$ on the successive
ranges $L=1$, $2\le L\le3$, $4\le L\le12$, $L\ge13$; each is at most
$L$. Monochromatic cycles contribute zero. Sum the run lengths.
\end{outline}

\begin{theorem}[Theorem C: the universal two-class ceiling]\label{thm:universal-capacity}
\formal{CapacityUniversal.lean}{IsSharedCode.two_class_capacity_bound}
For every field, every $d,b\ge1$, every partition into at most two classes,
and every linear source with deterministic shared decoders,
\begin{equation}\label{eq:universal-capacity}
 D\le\frac{41}{64}bd^2+\frac53 b d^{3/2}.
\end{equation}
In particular the normalized dimension is at most $41/64+o(1)$,
uniformly in the partition, field, word width and decoders.
\end{theorem}
\begin{outline}
Apply \cref{thm:iterated-tp} with successive shifts $s,3s,9s$.
The three difference sets contain every multiple of $s$ with coefficient
in $[-1,1]$, $[-4,4]$, and $[-13,13]$, respectively: this follows by
choosing difference digits from $\{-1,0,1\}$ in bases $1,3,9$.
Reduction modulo $d$ may identify elements, which causes no problem.
The component estimates give
\[
 D\le\frac12bd^2+\frac14bdQ_s+
 \frac18bd\bigl(R_4(s)+R_{13}(s)\bigr)+\frac34bdg_s.
\]
Use \eqref{eq:run-budget} and \eqref{eq:doubled-runs} to obtain
\begin{equation}\label{eq:one-step-universal}
 D\le\frac9{16}bd^2+\frac3{16}bdQ_s-\frac1{32}bdQ_{2s}
       +\frac34bdg_s.
\end{equation}
The negative term is retained before averaging.

Choose $a$ uniformly from all $d$ residues and put
$A_r=\mathbb E_a Q_{2^ra}$ and $B_r=\mathbb E_a\gcd(d,2^ra)$.
If the label sizes are $n_0,n_1$, ordered-pair counting gives
$A_0=2n_0n_1/d\le d/2$. Average \eqref{eq:one-step-universal} at
$s=2^ra$, then weight these inequalities by $w_r=(5/6)6^{-r}$ for
$r\ge0$. The weights sum to one and all terms are bounded for fixed
$d$. Since $(3/16)w_r=(1/32)w_{r-1}$ for $r\ge1$, the internal
$A_r$ terms cancel, leaving
\[
 D\le\frac{41}{64}bd^2+\frac34bd\sum_{r\ge0}w_r B_r.
\]
Equivalently, sum through a finite index and then let the index grow;
the terminal term has absolute value at most $bd^2w_r/32$ and tends to zero.
No uniformity of $2^ra$ on the whole cycle is assumed.

Let $a_r=\gcd(d,2^r)$ and $n=d/a_r$. Its actual distribution is uniform
on the multiples of $a_r$, so
$B_r=a_r\mathbb E_{u\in\Z/n}\gcd(n,u)$.
For each divisor $e$ of $n$, at most $n/e$ residues have gcd equal to
$e$. Thus the mean gcd is at most the number of divisors of $n$,
which is at most $2\sqrt n$ by pairing divisors across $\sqrt n$.
This includes $u=0$ and $n=1$. Consequently
\[
 B_r\le2\sqrt d\,2^{r/2},\qquad
 \sum_{r\ge0}w_rB_r\le
 \frac{5\sqrt d}{3(1-\sqrt2/6)}\le\frac{20}{9}\sqrt d.
\]
Substitution proves \eqref{eq:universal-capacity}. The proof works with
rows in $(F^b)^d$ throughout; no lane decomposition of $C$ was used.
\end{outline}

The bound improves the general two-class $3/4+o(1)$ ceiling for linear
decoders stated in \cite[Open problems]{SeventeenthPaper2026}, and allows
arbitrary deterministic decoders. It leaves the conjectured
$1/2+o(1)$ ceiling open. The parameter being bounded is the dimension of
a linear correctness source, not the density of native algorithms.

\subsection{Arbitrarily many classes}

\begin{theorem}[Theorem M: a square-root class deficit]\label{thm:many-classes}
\formal{CapacityDifferenceCover.lean}{IsSharedCode.difference_cover_real_bound, IsSharedCode.square_cover_bound, IsSharedCode.saturated_cover_bound}
For positive integers $u,v$, set
$R=\min(uv,d-1)$ and $U=\min(d,M\lfloor d/(R+1)\rfloor)$.
Then
\begin{equation}\label{eq:many-classes}
 D\le N-\frac{bd(d-U)}{u+v}.
\end{equation}
If $M\le d/2$ and $t=\lceil\sqrt{2M}\rceil$, then
\begin{equation}\label{eq:sqrt-classes}
 D\le\left(1-\frac1{4t}\right)N.
\end{equation}
For every $d\ge2$ and every $M$, writing $t_d=\lceil\sqrt{d-1}\rceil$,
$D\le N-bd(d-M)/(2t_d)$.
\end{theorem}
\begin{outline}
Use the shift set
$E=\{0,1,\ldots,u-1\}\cup\{u,2u,\ldots,vu\}$, reduced modulo $d$.
Its size is at most $u+v$. Every integer $1\le h\le uv$ equals $ju-i$
with $j=\lceil h/u\rceil\le v$ and $0\le i<u$. Thus $E-E$ contains
every signed difference through $R$.

For one class, list its vertices cyclically. The positive gaps sum to
$d$, including the sole gap $d$ for a singleton. Keep consecutive edges
whose gap is at most $R$. If there are $\ell\ge1$ larger gaps, there
are at most $\ell$ components and $\ell(R+1)\le d$. If there is no
large gap, there is one component, also at most $\lfloor d/(R+1)\rfloor$
because $R\le d-1$. Hence $\chi(E)\le U$.
The translate bound gives $D\le N-bd(d-\chi(E))/|E|$.
Both numerators are nonnegative; replacing $\chi(E)$ by $U$ and
$|E|$ by its upper bound proves \eqref{eq:many-classes}.

For $u=v=t$, either $t^2\ge d-1$, giving $U=M\le d/2$, or
$U\le Md/(t^2+1)<d/2$. Equation \eqref{eq:sqrt-classes} follows.
For $u=v=t_d$ the radius is $d-1$ and $U=M$, giving the last claim.
Reduction of the shift set may create repetitions; the proof uses only
its upper cardinality bound. It includes $d=1$ in the first inequality.
\end{outline}

This application uses the elementary sparse difference cover above, not a
new difference-cover theorem. It changes the general deficit order from
$1/M$ to $1/\sqrt M$, while stronger fixed-$M$ bounds remain available.
For example, if $D\ge(1-\eps)N$ and $0<\eps\le1/8$, then
\eqref{eq:sqrt-classes} gives the precise dichotomy
\[
 M>d/2\quad\text{or}\quad M>1/(128\eps^2).
\]
Indeed $\lceil\sqrt{2M}\rceil\ge1/(4\eps)$ forces
$\sqrt{2M}>1/(4\eps)-1\ge1/(8\eps)$.
The saturation at $d$ is necessary because singleton classes can decode
the full source. Nothing here extracts these decoder interfaces from a
native tape execution or closes them under composition.
